% 編譯：cd docs && xelatex 戰術語言.tex
\documentclass[11pt]{article}
\usepackage[a4paper,margin=2cm]{geometry}
\usepackage{fontspec}
\usepackage{xeCJK}
\usepackage{booktabs}
\usepackage{enumitem}
\usepackage{hyperref}

\setmainfont{Times New Roman}
\IfFontExistsTF{Songti TC}{
  \setCJKmainfont{Songti TC}
  \setCJKsansfont{Heiti TC}
}{
  \setCJKmainfont{FandolSong}
  \setCJKsansfont{FandolHei}
}

\setlist{nosep,leftmargin=1.8em}
\pagestyle{empty}

\newcommand{\term}[1]{\textbf{#1}}

\begin{document}

{\LARGE\bfseries 戰術語言一覽\par}

\section{傳球相關}

\subsection*{\term{傳球}}
球由傳球者的腳邊移動至接球者的腳邊，雙方皆須觸碰到球。不符合下列任何特殊傳球時，一律標為傳球。

\subsection*{\term{直塞}（Line breaking pass）}
傳球向前且穿越防線。
\begin{itemize}
  \item \textbf{起點}：任意區域。
  \item \textbf{終點}：比起點更靠近球門，且位於 \term{弧頂前緣} 或 \term{禁區中央}。
  \item \textbf{特徵}：終點深度穿過當下防線。
\end{itemize}

\subsection*{\term{回做球}}
球由前方區域往回傳給身後隊友。
\begin{itemize}
  \item \textbf{常見路徑}：弧頂前緣 $\rightarrow$ 左肋深處／右肋深處。
  \item \textbf{特徵}：終點 zone 深度小於起點（往回傳）。
\end{itemize}

\subsection*{\term{盤帶推進}}
同一接球者連續兩次觸球，球與人同步推向球門。
\begin{itemize}
  \item \textbf{執法}：兩拍為同一人，且後一拍比前一拍更靠近球門。
\end{itemize}

\subsection*{\term{強弱邊轉移}（Switch of play）}
長傳橫跨半場，利用弱側空間。
\begin{itemize}
  \item \textbf{常見路徑}：右邊路／右路中路 $\rightarrow$ 左邊路／左路底線；或後場中路長傳至左路底線。
  \item \textbf{特徵}：橫跨半場換邊，或自中低位縱深大跳躍換邊。
\end{itemize}

\subsection*{\term{倒三角傳球}（Cut-back）}
自底線或禁區深處貼地回傳弧頂前緣。
\begin{itemize}
  \item \textbf{起點}：左路底線、右路底線，或禁區中央深處。
  \item \textbf{終點}：弧頂前緣。
\end{itemize}

\subsection*{\term{傳中}（Cross）}
邊路傳球至禁區中央。
\begin{itemize}
  \item \textbf{起點}：左路底線或右路底線（或邊路前段）。
  \item \textbf{終點}：禁區中央。
\end{itemize}

\subsection*{\term{過頂長傳}（Lobbed pass / Over the top）}
中低位起腳，縱深大跳躍越過防線。
\begin{itemize}
  \item \textbf{常見路徑}：後場中路 $\rightarrow$ 左路底線；或中低位 $\rightarrow$ 弧頂／禁區／底線，縱深跳躍 $\geq 2$ 格。
  \item \textbf{特徵}：球路越過防線頭頂至身後空間。
\end{itemize}

\subsection*{\term{第三人出球}（破高位壓迫）}
在逼搶壓力下，第二人短傳給已拉開的第三人脫困。
\begin{itemize}
  \item \textbf{常見序列}：後場禁區 $\rightarrow$ 後場中路 $\rightarrow$ 邊路後段／邊後衛，再接肋部或前場推進。
  \item \textbf{特徵}：第二、第三人接力，避開第一線壓迫。
\end{itemize}

\section{終點結果}

\subsection*{\term{射正}}
終點在禁區中央，球進入球框範圍（未進特定得分區）。

\subsection*{\term{得分}}
終點在禁區中央，球進入球框之特定得分範圍。

\section{戰術區 20 格名稱}

完整示意見戰術板。編號與名稱對照：

\vspace{0.3em}
\small
\begin{tabular}{@{}clcl@{}}
\toprule
\textbf{Zone} & \textbf{名稱} & \textbf{Zone} & \textbf{名稱} \\
\midrule
1 & 後場左角 & 11 & 左邊路 \\
2 & 後場禁區 & 12 & 左邊路（前段） \\
3 & 後場右角 & 13 & 左肋 \\
4 & 左邊路（後段） & 14 & 弧頂前緣 \\
5 & 左邊路（中段） & 15 & 右肋 \\
6 & 左肋深處 & 16 & 右路中路 \\
7 & 後場中路 & 17 & 右邊路 \\
8 & 右肋深處 & 18 & 左路底線 \\
9 & 右邊路（後段） & 19 & 右路底線 \\
10 & 右邊路（中段） & 20 & 禁區中央 \\
\bottomrule
\end{tabular}
\normalsize

\section{戰術描述常用場地名稱}

戰術文字常優先使用下列簡稱（皆對應上表某一 zone）：

\vspace{0.4em}
\begin{tabular}{@{}ll@{}}
\toprule
\textbf{名稱} & \textbf{說明} \\
\midrule
後場中路 & 後場中央帶 \\
左肋深處／右肋深處 & 後場至中場的肋部深區 \\
左邊路／右邊路 & 邊路中段 \\
左肋／右肋 & 進攻肋部 \\
弧頂前緣 & 禁區前沿弧頂 \\
右路中路 & 右側中場通道 \\
左路底線／右路底線 & 底線角區 \\
禁區中央 & 大禁區中央 \\
\bottomrule
\end{tabular}


\end{document}
